From \parencite{peffers2007design_science_research_methodology}:

Om DSR generelt
"Design science creates and evaluates IT artifacts intended to solve identified organizational problems."
"It involves a rigorous process to design artifacts to solve observed problems, to make research contributions, to evaluate the designs, and to communicate the results to appropriate audiences."

Activity 1 og 2: Problem og objectives
"Define the specific research problem and justify the value of a solution."
"Infer the objectives of a solution from the problem definition and knowledge of what is possible and feasible."

Activity 3: Design and development
"Create the artifact... This activity includes determining the artifact's desired functionality and its architecture and then creating the actual artifact."

Activity 4: Demonstration
"Demonstrate the use of the artifact to solve one or more instances of the problem. This could involve its use in experimentation, simulation, case study, proof, or other appropriate activity."

Activity 5: Evaluation
"Observe and measure how well the artifact supports a solution to the problem. This activity involves comparing the objectives of a solution to actual observed results from use of the artifact in the demonstration."

Om ikke-sekvensiell bruk
"This process is structured in a nominally sequential order; however, there is no expectation that researchers would always proceed in sequential order from activity 1 through activity 6."

Om design-centered inngang, som passer min situasjon
"A design- and development-centered approach would start with activity 3. It would result from the existence of an artifact that has not yet been formally thought through as a solution for the explicit problem domain in which it will be used."

Iterations:
At the end of this activity (evaluation) the researchers can decide whether to iterate back to activity 3 to try to improve the effectiveness of the artifact or to continue on to communication."

From \parencite{hevner2004design_science_information_systems_research}

Hva DSR er (kjernedefinisjonen)
"The design-science paradigm seeks to extend the boundaries of human and organizational capabilities by creating new and innovative artifacts."
"Design science creates and evaluates IT artifacts intended to solve identified organizational problems."

Hva et artefakt er
"IT artifacts are broadly defined as constructs (vocabulary and symbols), models (abstractions and representations), methods (algorithms and practices), and instantiations (implemented and prototype systems)."

Om kunnskapsbygging gjennom design
"In the design-science paradigm, knowledge and understanding of a problem domain and its solution are achieved in the building and application of the designed artifact."